\section{Discussion and limitations}
\subsection{The joint reading of prices and quantities}

The central empirical pattern has three parts. Event-window addition returns are substantially lower in the latest block than in the unusually high first block. A positive average remains visible when the later decade is pooled. At the same time, trading becomes more concentrated immediately before implementation and covered benchmark-linked funds change their holdings in the index-directed direction. These facts can coexist because trading quantities, price movements and the timing of price adjustment are distinct objects.

A market with more predictable benchmark demand and more responsive liquidity supply could display this combination. So could a market in which part of the measured return reflects selection and information while predictable trading is executed with limited additional price movement. The current evidence does not distinguish these accounts sharply enough to rank their causal importance. The appropriate inference is a constraint on simple stories: low implementation-window returns do not establish negligible index-related activity, while large holdings changes do not establish a large causal price response.

The comparison with \citet{greenwood2025} is therefore most useful as a common empirical question. Sweden differs in index rules, market composition, fund coverage and the observed historical interval. The Swedish early sample is small and does not extend to the decades with the largest US estimates. A numerical difference between countries cannot be attributed to passive ownership, liquidity or arbitrage capacity without harmonized measurement and a design for that comparison. This study contributes evidence on a different setting rather than a rejection or confirmation of a universal index-effect trajectory.

The ratio evidence illustrates a broader measurement issue. If a historical ownership database captures a growing fraction of the relevant investors, its measured demand can increase even without a commensurate increase in total demand. A falling return-to-measured-demand ratio can then reflect coverage as well as price responsiveness. This is not a reason to abandon holdings data. It is a reason to report common samples, mandate definitions, report dates, denominator uncertainty and independent reporting comparisons alongside the ratio itself.

\subsection{Threats to interpretation}

The first limitation is identification. Reconstitution is related to size, trading and industry composition. Market adjustment and observed controls do not remove all reasons that selected securities might have different returns. The current comparison-security construction is not a reconstructed threshold experiment, and the share-class exercise can contain spillovers. Consequently, the paper estimates return patterns around events and conditional associations, rather than treatment effects of index assignment or structural stock-demand elasticities.

The second limitation is historical measurement. The event inventory is larger than the usable return panel, and a portion of the notice URLs could not be recovered during the independent retrieval audit. A successfully retrieved page is not by itself proof that every event field has been independently reconciled. Fund holdings depend on report completeness, mandate classification and historical share-count adjustments. Checksums and code replication secure the identity and transformation of the supplied files; they do not authenticate the economic content of each vendor observation.

The third limitation is statistical resolution. The final block has only nine reviews per event type, early holdings estimates have still fewer, and events can reappear across reviews. Alternative dependence estimators and review influence checks help reveal sensitivity but do not provide exact finite-sample coverage. The latest return intervals are wide relative to economically plausible effects. An insignificant coefficient is therefore a weak basis for saying that an effect is absent, particularly when equivalence is not established.

The fourth limitation is the history of analytical choices. The repository contains several specifications, screens and previously explored breakpoints. The new audit has access to those findings. This study cannot claim that its full hypothesis family, periodization or extension was fixed independently of outcomes. The ten-test Holm calculation is a bounded sensitivity exercise, not a comprehensive correction for that history. The longer-horizon result is openly exploratory and should motivate future validation rather than be presented as a pre-specified discovery.

Finally, the code audit identifies numerical rather than purely economic uncertainty in some auxiliary intervals. Re-estimation reproduces the point estimates and sample counts, but does not reproduce every bootstrap endpoint identically. The largest discrepancies occur in sparse auxiliary ratio designs. The manuscript avoids making substantive claims from those endpoints and reports analytical and Fieller alternatives with their own assumptions. This separation matters: successful execution of a pipeline is weaker than validation of every inferential procedure it implements.

\subsection{A focused route to stronger evidence}

The most valuable next step is to improve the link between the observed event and a credible counterfactual, rather than to add more significance tests. A historical reconstruction of the Nasdaq ranking and eligibility rules could identify whether a defensible local comparison exists. Such a design would require the actual selection inputs at each review, the relevant industry constraints and buffers, an analysis of manipulation or sorting, and evidence that observations near any boundary are comparable. A market-capitalization screen alone is insufficient.

A second priority is to reconcile the holdings measurement at the report and share-class level. The decisive records are complete pre- and post-event reports, dated mandates, corporate-action adjustments and the exact shares-outstanding denominator. A stable-coverage panel would permit a more credible comparison of measured demand over time. Separate analyses using strict followers and broader benchmark-linked investors could then quantify the sensitivity to the investor definition without conflating it with changing event samples.

A third priority is prospective validation. Future OMXSB reviews could be analyzed under a dated protocol specifying the sample, windows, benchmark, exclusion rules and primary tests before their returns are observed. The current 2026 block is incomplete, so further events are natural additions to a new validation sample. That would test whether the later-decade patterns persist without retroactively treating the existing exploratory results as confirmatory. Until then, the present study's value rests on transparent measurement, reproducible arithmetic and appropriately bounded conclusions.
